\chapter{Kernel-Level Characterization}
\label{ch:characterization}

This chapter presents the classification of cuVSLAM's 49 GPU kernels across the 27-sequence corpus and subjects that classification to four independent validations: input stability (\S\ref{sec:char-stability}), unsupervised recovery by two clustering algorithms (\S\ref{sec:char-clustering}), cross-microarchitecture agreement (\S\ref{sec:char-crossdev}), and a known-truth calibration suite (\S\ref{sec:char-calibration}). It closes with measurement of host--device transfers at the system edge (\S\ref{sec:char-transfers}).

\section{What the Screen Reveals}

The nsys timeline over full sequences shows the expected structure of a frame-synchronous perception pipeline: a small set of front-end kernels launched every frame (image cast, Gaussian pyramid scaling, gradient convolutions, Lucas--Kanade tracking, and the matcher pair) accounts for the bulk of launches ($\sim$80 per frame; $\sim$200{,}000 launches over a 2{,}500-frame sequence). The bundle-adjustment solver family (\code{sba::*}, plus cuSOLVER helper kernels) forms a second tier, while SLAM-layer keyframe kernels (\stBuild{}, \stTrack{}) launch sparsely per keyframe and per loop-closure scan, respectively. The NVTX stage table (measured, not name-matched; \S\ref{sec:bg-tools}) confirms stage membership: \stTrack{} executes inside the loop-closure-and-optimization range, \stBuild{} within keyframe ingestion the measured anchor for every loop-closure claim that follows.

\section{Class Stability Across Inputs}
\label{sec:char-stability}

The classification model presented in \S\ref{sec:met-tree} was run independently on each sequence, and then aggregated over the entire dataset. The major conclusion is that there is \textbf{91\% of agreement in the modal class}: out of 49 kernels, 24 belong to the exact same bottleneck class for all 27 sequences, while 42 belong to the same class in at least 80\% of them. The bottleneck class is, to the first order, \emph{the property of the kernel, not the input}.

The residual flips are not noise; they are physics. They cluster around kernels whose working set sits near the 25\,MB L2 capacity: on small indoor sequences the set fits (L2-reuse behaviour, G3), on street-scale it spills (bandwidth behaviour, G1). The sensitivity analysis (\S\ref{sec:met-sensitivity}) flags exactly these kernels as borderline, and the substrate synthesis (Chapter~\ref{ch:synthesis}) treats their verdicts as workload-conditional 25 of 49 verdicts flip with workload scale, always in the physically expected direction.

\section{The Taxonomy Is Discovered, Not Asserted}
\label{sec:char-clustering}

The decision tree assigns labels; whether the \emph{classes themselves} are real structure in the data is tested without labels. Figure~\ref{fig:taxonomy} first shows how the 49 characterized kernels distribute across the classes. Pooling every kernel's feature vector (memory/compute/DRAM speed-of-light, LFMR, occupancy, sectors-per-request, and the two dominant stall shares) across sequences and running k-means over a range of cluster counts, the silhouette criterion selects \textbf{$k = 7$--$8$} matching the seven populated classes plus the screened remainder, with cluster purity 0.68 and adjusted Rand index 0.30 against the tree labels. Using a completely different algorithm Ward hierarchical clustering~\cite{ward1963}, cut at $k{=}8$ the same agreement is reproduced: \textbf{ARI 0.31, purity 0.675}. Two unrelated algorithms, given no thresholds and no labels, find the same $\sim$8-way structure encoded by the tree. The tree should be read as the \emph{labelling} of natural structure and the clustering as its \emph{independent validation} precisely DAMOV's discipline, reproduced on GPU data.

\begin{figure}[t]
\centering
\begin{tikzpicture}
\begin{axis}[
  width=0.82\textwidth, height=5cm,
  ybar, bar width=17pt,
  ymin=0, ymax=15, ylabel={kernels (of 49)},
  symbolic x coords={G0,G1,G2,G3,G4,G5,G6,G7},
  xtick=data, tick label style={font=\small}, ylabel style={font=\small},
  nodes near coords, nodes near coords style={font=\scriptsize},
  enlarge x limits=0.08,
]
\addplot[draw=black!40,fill=cPIM!45] coordinates
  {(G0,1)(G1,9)(G2,10)(G3,9)(G4,5)(G5,2)(G6,0)(G7,13)};
\end{axis}
\end{tikzpicture}
\vspace{-2mm}
\caption[Kernel distribution across the taxonomy.]{Distribution of the 49
characterized kernels across the eight classes (campaign modal class). The
near-data-affine classes (G1, G2, G4; orange leaves in Figure~\ref{fig:tree})
account for 24 kernels. G6 is unpopulated in cuVSLAM (no kernel is on-chip-bound
with DRAM idle), and G7, the class that emerged from the data, is the single
largest, capturing the dependency-stalled tracking and solver kernels.}
\label{fig:taxonomy}
\end{figure}

\section{Cross-Microarchitecture Agreement}
\label{sec:char-crossdev}

If the classes captured artifacts of one GPU, they would not persist across microarchitectures. The same workload (TUM office) was characterized independently on two deliberately different devices: the \code{sm\_89} workstation and an \code{sm\_75} laptop with a quarter the memory and a third the power budget. \textbf{36 of 45 signal kernels (80\%) receive the same class on both} (76.6\% including the launch-tax G0 tier). The honest caveat: the two devices differ in memory system as well as core generation, so some disagreements are physically real L2-capacity crossovers rather than classification error making 80\% a conservative floor for device-independence of the classification. This mirrors DAMOV's core-type-independence check, under a stricter test.

\section{The Calibration Suite: Classifying Known Truth}
\label{sec:char-calibration}

DAMOV validated frozen thresholds on held-out functions with known behaviour. Here, a \emph{calibration suite} of eight archetype kernels is used, each with a known ground-truth class by construction: a streaming triad (G1), a random gather (G2), an L2-resident sweep (G3), a coalesced pointer-chase (G4), an FMA polynomial (G5), a bank-conflicted shared-memory kernel (G6), a single-warp dependency chain (G7), and a sub-screen tiny kernel (G0). These are compiled and pushed through the identical capture-and-classify pipeline, blind. The first pass classified five of eight as designed; analysis of the three misses showed each to be a defect in the \emph{archetype}, not the classifier: a per-lane-random pointer-chase is both scattered and latency-bound (G2+G4); a barrier-heavy shared-memory kernel presents as dependency-stalled rather than G6; and a sub-floor kernel is legitimately \emph{screened}, as per the pipeline's Step-1 semantics. Each archetype was corrected to embody a single class \emph{with the classifier untouched throughout} and the suite is retained in the artifact as a permanent regression harness: any future change to thresholds or metrics must re-classify known truth correctly.

\section{Near-Data Affinity of the Workload}
\label{sec:char-affinity}

Rolling the per-kernel classes up by measured GPU time: across the 27-sequence campaign, \textbf{21\% of GPU time is in strong-affinity classes, 40\% in conditional-affinity classes, and 28\% in no-affinity classes} (the remainder are sub-floor). Thus, roughly \textbf{61\% of GPU time carries near-data affinity} under the taxonomy's default mapping, rising to $\sim$72\% time-weighted in deep single-sequence studies where the loop-closure path is exercised densely. This is the quantity the placement model in Chapter~\ref{ch:synthesis} prices. It is reported here with its composition, not as a headline slogan, since the conditional 40\% is exactly the portion whose verdict depends on the scatter-PiM-vs-layout-fix question addressed in the locality and attribution chapters.

\section{The System Edge: Host--Device Transfers}
\label{sec:char-transfers}

Independent of any kernel's internals, the timeline measures the pipeline's boundary traffic: explicit host-to-device and device-to-host copies account for \textbf{41\% of total kernel time} on the RGB-D workload, and the per-frame sensor upload is \textbf{1.68\,MB/frame} the camera image and depth crossing PCIe into GPU DRAM before any kernel touches them. This is the single most direct near-sensor argument in the thesis: it is traffic that exists only because the pixels are \emph{born} on the wrong side of the bus, and no on-GPU optimization can remove it. The host-side companion measurements (66\,MB of storage ingestion per run; 708\,MB peak host memory) are detailed in the synthesis chapter, where the boundary picture is assembled end-to-end.